% Generated draft: requires mathematical and source review.
\documentclass[11pt]{article}
\usepackage[margin=25mm]{geometry}
\usepackage{fontspec}
\setmainfont{Latin Modern Roman}
\usepackage{amsmath,amssymb,mathtools,unicode-math}
\setmathfont{Latin Modern Math}
\usepackage{enumitem,needspace,xurl,hyperref}
\setlength{\parindent}{0pt}
\setlength{\parskip}{5pt}
\setlength{\emergencystretch}{3em}
\newcommand{\N}{\mathbb{N}}
\newcommand{\Z}{\mathbb{Z}}
\newcommand{\Q}{\mathbb{Q}}
\newcommand{\R}{\mathbb{R}}
\newcommand{\C}{\mathbb{C}}
\newcommand{\F}{\mathbb{F}}
\newcommand{\A}{\mathbb{A}}
\newcommand{\PP}{\mathbb{P}}
\newcommand{\E}{\mathbb{E}}
\newcommand{\eps}{\varepsilon}
\newcommand{\dd}{\,\mathrm{d}}
\begin{document}
\section*{UnsolvedMath ID: 30006519}
\textbf{Source problem code:} catalogue entry
\subsection{Definitions and mathematical statement}
Let \(C(0)\) denote the connected component containing the distinguished origin \(0\) in the random structure called the \emph{infinite-noodle model} in the supplied record, and let \(\mathbb P\) denote its probability law. Write \(X=|C(0)|\) for its size under the intended model's size convention, allowing \(X=\infty\). The supplied record does not define the model, its probability law, or what is counted by component size. Consequently, the following formulation retains those unresolved source-dependent conventions and is not a self-contained specification of the random object.

The two conjectural assertions to be established are
\[
\mathbb P\bigl(C(0)\text{ is finite}\bigr)=1
\]
and, with \(\alpha=(2\sqrt{2}-1)/7\),
\[
\mathbb P(X\geq k)=k^{-\alpha+o(1)}
\qquad\text{as }k\to\infty,\quad k\in\Z_{\geq1}.
\]
Precisely, the second assertion means that there exists a real-valued function \(r\) on the positive integers, with \(r(k)\to0\), such that
\[
\mathbb P(X\geq k)=k^{-\alpha+r(k)}
\]
for every sufficiently large positive integer \(k\). No assertion is made here about a multiplicative asymptotic constant.

The supplied original excerpt attributes almost-sure finiteness to its reference [4] and the tail conjecture to its reference [1], Conjecture 1.10. The corresponding numbered bibliography is absent, so these reference labels cannot be reliably assigned to the supplied literature links.
\subsection{Short English statement}
In the intended infinite-noodle model, is the component containing the origin finite almost surely, and does its size have tail probability \(k^{-(2\sqrt{2}-1)/7+o(1)}\) as positive integers \(k\) tend to infinity? The model and the convention for measuring component size require clarification from the cited source.
\subsection{Sources}
\begin{itemize}
\item Supplied local UnsolvedMath registry record, permanent ID \(30006519\), title: \emph{Finiteness and Tail Decay of Infinite-Noodle Components}; exact source problem code: \texttt{catalogue entry}. The record labels the problem open and supplies no named proposer.
\item \emph{Enumerative Combinatorics} (2026), Oberwolfach Reports, pp.~81--156, DOI \texttt{10.4171/owr/2026/2}. Supplied report provenance: \url{https://doi.org/10.4171/owr/2026/2}. The background describes Oberwolfach Reports as documenting mathematical workshops and their open-problem discussions.
\item F\'eray--Th\'evenin, \emph{Components in meandric systems and the infinite noodle}. Supplied literature reference: \url{https://arxiv.org/abs/2201.11572}.
\item Curien et al., \emph{Uniqueness of the infinite noodle}. Supplied literature reference: \url{https://arxiv.org/abs/1701.01083}.
\item Supplied literature assessment dated \texttt{2026-08-22}: the component of zero is not known to be finite almost surely, and the stated power-law tail conjecture remains unresolved. This is supplied dated curation derived from the OWR record, not an independently verified literature assessment. The source itself requests independent verification of time-sensitive status and cited claims.
\item The registry assigns the statement the label \texttt{corrected\_verified}, explaining that only formula delimiters, notation, grammar, and typography were repaired while preserving mathematical content and scope. This supplied editorial label does not establish independent mathematical verification.
\end{itemize}
\end{document}
